\chapter{假设检验、局部似然几何与 Pearson 投影}\label{chap:math-stat-testing}

估计问题在参数空间中寻找一个点或区域，检验问题则比较两个参数集合。有限样本精确检验依赖特殊模型中的枢轴量或离散分布；一般正则模型的大样本检验则由 \chapref{chap:math-stat-estimation-likelihood} 的 局部渐近正态 二次型统一控制。因而本章不按“某某检验公式”分类，而是先定义检验作为一个受第一类错误约束的决策问题，再从 Neyman--Pearson 的有限样本最优性进入 Fisher 几何，随后由同一个局部法向投影推出 LR、Wald、得分、局部功效与 干扰参数 消元，最后把 Pearson 卡方解释成概率单纯形上的同一投影机制。

\section{检验函数、水平、功效与 Neyman--Pearson 原理}

\begin{definition}[随机化检验]
设参数空间分成原假设与备择集合
\[
 \Theta_0\subseteq\Theta,
 \qquad
 \Theta_1=\Theta\setminus\Theta_0.
\]
一个随机化检验是可测函数
\[
 \phi_n:\mathcal X^n\to[0,1].
\]
观察到 $X^n=x^n$ 后，以概率 $\phi_n(x^n)$ 拒绝原假设。普通拒绝域检验对应 $\phi_n\in\{0,1\}$。
\end{definition}

\begin{definition}[水平、大小与功效]
检验的功效函数定义为
\[
 \operatorname{Power}_n(\theta)
 =\mathbb E_\theta\phi_n(X^n).
\]
在 $\theta\in\Theta_0$ 时，它就是第一类错误概率；在 $\Theta_1$ 上则表示识别备择的概率。若
\begin{equation*}\sup_{\theta\in\Theta_0}\mathbb E_\theta\phi_n\le\alpha,
\end{equation*}
则称检验的水平不超过 $\alpha$。因此检验的基本优化顺序是：先满足整个原假设上的错误控制，再在这个约束下比较功效。
\end{definition}

\begin{definition}[对复合零假设有效的 p 值]
$p$ 值也可以直接用错误控制定义。随机变量 $p_n(X^n)\in[0,1]$ 若满足
\[
 \sup_{\theta\in\Theta_0}
 \mathbb P_\theta\{p_n(X^n)\le u\}
 \le u,
 \qquad 0\le u\le1,
\]
则称其为对 $\Theta_0$ 有效的 $p$ 值。这样“当 $p_n\le\alpha$ 时拒绝”自动给出水平不超过 $\alpha$ 的检验。若一个统计量 $T$ 越大越反对原假设，一个常见构造是
\[
 p_n(x^n)
 =\sup_{\theta\in\Theta_0}
 \mathbb P_\theta\{T(X^n)\ge T(x^n)\}.
\]
简单原假设时它退化为普通尾概率；复合原假设中的上确界则保证 干扰参数 不能被随意代成有利于拒绝的值。
\end{definition}

\begin{theorem}[Neyman--Pearson 最强性]
先看有限样本的基本最优性问题：
\[
 H_0:\theta=\theta_0,
 \qquad
 H_1:\theta=\theta_1.
\]
Neyman--Pearson 引理说明，在所有满足给定水平约束的检验中，最强检验按 似然比
\[
 \frac{p_{\theta_1}^{(n)}(X^n)}{p_{\theta_0}^{(n)}(X^n)}
\]
排序样本点；超过阈值时拒绝，边界处必要时随机化。

设简单零假设与备择的密度 \(p_0,p_1\) 具有共同支配测度。选 \(c\ge0\)，令
\(\varphi^*(x)=1\) 当 \(p_1(x)>cp_0(x)\)，
\(\varphi^*(x)=0\) 当 \(p_1(x)<cp_0(x)\)，
等号处可随机化，使 \(\mathbb E_0\varphi^*=\alpha\)。
则对任何满足 \(\mathbb E_0\varphi\le\alpha\) 的检验 \(\varphi\)，有
\(\mathbb E_1\varphi\le\mathbb E_1\varphi^*\)。
\end{theorem}
\begin{proof}
在 \(p_1-cp_0>0\) 处，\(\varphi^*=1\ge\varphi\)；在
\(p_1-cp_0<0\) 处，\(\varphi^*=0\le\varphi\)。
因此逐点有
\((\varphi^*-\varphi)(p_1-cp_0)\ge0\)。
积分后得
\[
\mathbb E_1(\varphi^*-\varphi)
\ge c\mathbb E_0(\varphi^*-\varphi)
\ge0,
\]
末步使用检验大小约束。等号集合的随机化只为精确达到指定大小。

把逐点不等式按拒绝域分段积分，可直接看出水平约束为何给出功效最优。

令 $\phi^*$ 为 似然-比 检验，阈值为 $c\ge0$，拒绝域为
\[
 R=\{t:p_{\theta_1}^{(n)}(t)>c\,p_{\theta_0}^{(n)}(t)\},
\]
即 $\phi^*(t)=\mathds{1}_R(t)+\gamma(t)\mathds{1}_{\{p_1=cp_0\}}(t)$，其中 $0\le\gamma(t)\le1$ 为边界随机化。令 $\phi$ 是另一个满足 $0\le\phi\le1$ 且
\[
 \mathbb E_{\theta_0}\phi\le\mathbb E_{\theta_0}\phi^*
\]
的检验。把样本空间按 $R$ 与 $R^c$ 分段讨论。在 $R$ 上，由定义 $p_{\theta_1}^{(n)}-c\,p_{\theta_0}^{(n)}>0$，且 $\phi^*=1$，而 $0\le\phi\le1$ 给出 $\phi^*-\phi\ge0$；在密度差严格为负的部分，$\phi^*=0$，给出 $\phi^*-\phi\le0$；密度差为零时，无论边界如何随机化，乘积都为零。因此所有部分中 $(\phi^*-\phi)$ 与 $(p_{\theta_1}^{(n)}-c\,p_{\theta_0}^{(n)})$ 同号，因此逐点有
\[
 (\phi^*-\phi)
 (p_{\theta_1}^{(n)}-cp_{\theta_0}^{(n)})\ge0.
\]
这一逐点同号性是 Neyman--Pearson 的几何核心：拒绝域 $R$ 选在 似然-比 大的方向，把检验函数差与密度差对齐。对样本空间积分，展开左端得
\begin{align*}
0&\le\int(\phi^*-\phi)(p_{\theta_1}^{(n)}-cp_{\theta_0}^{(n)})\,\dd\mu\\
&=\int(\phi^*-\phi)p_{\theta_1}^{(n)}\,\dd\mu
 -c\int(\phi^*-\phi)p_{\theta_0}^{(n)}\,\dd\mu\\
&=\mathbb E_{\theta_1}(\phi^*-\phi)-c\,\mathbb E_{\theta_0}(\phi^*-\phi).
\end{align*}
移项得到
\[
 \mathbb E_{\theta_1}(\phi^*-\phi)
 \ge c\,\mathbb E_{\theta_0}(\phi^*-\phi).
\]
由水平约束 $\mathbb E_{\theta_0}\phi\le\mathbb E_{\theta_0}\phi^*$，右端因子 $\mathbb E_{\theta_0}(\phi^*-\phi)\ge0$，又 $c\ge0$，故
\[
 \mathbb E_{\theta_1}(\phi^*-\phi)\ge0,
\]
即 $\mathbb E_{\theta_1}\phi^*\ge\mathbb E_{\theta_1}\phi$。因此 $\phi^*$ 在 $\theta_1$ 下的功效不低于任意同水平检验。注意这个结论只比较一个具体备择点；复合备择需要额外结构。
\end{proof}

\begin{theorem}[单调似然比与一致最强检验]
设统计量 $T$ 充分，具有共同支撑的正密度单调似然比族；阈值及随机化选在边界参数处达到大小 $\alpha$。
若一参数模型对统计量 $T$ 具有单调 似然比，即 $\theta_1>\theta_0$ 时
\[
 \frac{f_{\theta_1}(t)}{f_{\theta_0}(t)}
\]
关于 $t$ 单调不减，则 Neyman--Pearson 的拒绝域可以统一为 $T$ 的右尾。更精确地，若一个水平不超过 $\alpha$ 的检验 $\phi^*$ 满足：对任意其他同水平检验 $\phi$ 与每个 $\theta\in\Theta_1$，都有
\[
 \mathbb E_\theta\phi^*\ge\mathbb E_\theta\phi,
\]
则称 $\phi^*$ 在该复合备择上 一致最强（UMP）。Karlin--Rubin 定理说明：在一参数单调 似然-比 族中，检验
\[
 H_0:\theta\le\theta_0,
 \qquad
 H_1:\theta>\theta_0
\]
可以由 $T$ 的同一个右尾拒绝域得到 一致最强检验。双侧复合问题通常没有 一致最强检验，这不是技术缺陷，而是不同备择方向对“最有利的拒绝域”存在冲突。
\end{theorem}

\begin{proof}
设 $T$ 的分布族具有单调 似然比，并取阈值 $c$ 及必要的边界随机化，使
\[
 \mathbb E_{\theta_0}\phi^*=\alpha,
 \qquad
 \phi^*(T)=\mathds{1}\{T>c\}+\gamma\mathds{1}\{T=c\}.
\]
对任意固定 $\theta_1>\theta_0$，比值
\[
 f_{\theta_1}(t)/f_{\theta_0}(t)
\]
随 $t$ 单调，因此 Neyman--Pearson 引理给出的最强 level-$\alpha$ 检验正是同一个右尾 $\phi^*$；关键在于拒绝域不随所选的 $\theta_1$ 改变。这一步把"对每个备择最强"压缩为"对每个备择都对应同一检验"，是 MLR 减少复杂性的核心机制。还要验证整个复合原假设上的水平。固定 $\theta<\theta_0$，令
\[
 r(t)=\frac{f_{\theta_0}(t)}{f_\theta(t)}.
\]
单调 似然比 使 $r(t)$ 单调不减，而右尾检验函数 $\phi^*(t)$ 也单调不减。在 $P_\theta$ 下，两个同向单调函数的协方差非负；用独立副本 $T'$ 可写成
\[
 2\Cov_\theta\{\phi^*(T),r(T)\}
 =\mathbb E_\theta\!\left[
 \{\phi^*(T)-\phi^*(T')\}
 \{r(T)-r(T')\}
 \right]\ge0,
\]
因为当 $T>T'$ 时两差同号，当 $T<T'$ 时两差同号反号，乘积恒非负。又因为 $\mathbb E_\theta r(T)=\int f_{\theta_0}=1$，所以
\[
 \mathbb E_{\theta_0}\phi^*(T)
 =\int\phi^*(t)f_{\theta_0}(t)\,\dd t
 =\int\phi^*(t)r(t)f_\theta(t)\,\dd t
 =\mathbb E_\theta\{\phi^*(T)r(T)\}
 \ge\mathbb E_\theta\phi^*(T)\cdot\mathbb E_\theta r(T)
 =\mathbb E_\theta\phi^*(T).
\]
中间的不等式即 $\Cov\ge0$。故对任意 $\theta\le\theta_0$，
\[
 \mathbb E_\theta\phi^*
 \le\mathbb E_{\theta_0}\phi^*=\alpha.
\]
因此它对整个复合原假设保持水平约束，又同时对每个 $\theta_1>\theta_0$ 都是最强检验，故为 一致最优势检验。这一论证也说明：若不同备择点的 似然-比 排序方向不能共用同一拒绝域，就不能仅靠 Neyman--Pearson 得到 一致最优势检验。
\end{proof}


\begin{exbox}[精确二项单侧检验]
若 $X_i\iid\mathrm{Bernoulli}(p)$，则
\[
 S=\sum_{i=1}^nX_i\sim\mathrm{Bin}(n,p)
\]
是充分统计量，并且 似然比 随 $S$ 单调。检验
\[
 H_0:p=p_0,
 \qquad
 H_1:p>p_0
\]
时，取最小整数 $c$ 使
\[
 \mathbb P_{p_0}(S\ge c)\le\alpha
\]
并在 $S\ge c$ 时拒绝，就得到有限样本水平不超过 $\alpha$ 的检验。由于二项分布离散，若要求 水平 恰好等于 $\alpha$，通常必须在边界随机化。该构造直接使用精确 二项 尾概率，因此结论对有限样本成立；若把二项尾替换为正态尾，则只得到大样本近似。
\end{exbox}

\begin{note}[双侧最优性与离散随机化]
\begin{workedexample}[一次试验的随机化阈值]
只观察一次 Bernoulli\((p)\) 试验，检验 \(p=p_0\) 对 \(p=p_1>p_0\)。当 \(0<\alpha<p_0\) 时，求在 \(X=1\) 上的随机化概率与检验功效。
\textbf{解。} 似然比在 \(X=1\) 时较大、在 \(X=0\) 时较小，因此最强检验只在 \(X=1\) 时以概率 \(\gamma\) 拒绝。其零假设拒绝概率为 \(p_0\gamma=\alpha\)，故 \(\gamma=\alpha/p_0\)；在 \(p_1\) 下功效为 \(p_1\gamma=\alpha p_1/p_0\)。
\end{workedexample}

双边备择通常不能共享同一个最强拒绝域：右侧备择要求右尾检验，左侧备择要求左尾检验。因此 Karlin--Rubin 的结论本质上是单边结论；双边问题必须改在无偏、相似或不变检验等较小类别中讨论最优性。

在离散模型中，非随机拒绝域未必能恰好达到 水平 $\alpha$。此时在 似然-比 边界上以概率 $\gamma\in[0,1]$ 拒绝，并选取 $\gamma$ 使
\[
\mathbb E_{\theta_0}\phi^*=\mathbb P_{\theta_0}(T>c)+\gamma\mathbb P_{\theta_0}(T=c)=\alpha.
\]
同一逐点不等式仍成立，因此随机化版本保持最强性。这里最强性仍只针对指定的简单备择；把它提升到复合备择正是后面需要 MLR 或其他结构的原因。
\end{note}

\begin{theorem}[有限群不变性下的精确检验]
Neyman--Pearson 从已知概率模型中寻找最优拒绝域；另一类有限样本精确检验并不需要把原假设分布完全参数化，而是利用原假设下的\emph{群不变性}。设有限群 $\mathcal G$ 作用于样本空间，$g x$ 表示变换后的样本，并假设在原假设下
\begin{equation}\label{eq:math-stat-group-invariance-null}
 X\overset{d}=gX,
 \qquad
 \text{对每个 }g\in\mathcal G.
\end{equation}
给定一个“越大越反对原假设”的统计量 $T$，定义轨道
\[
 \mathcal O(x)=\{gx:g\in\mathcal G\}
\]
以及随机化尾概率
\begin{equation*}p_{\mathcal G}(x)
 =\frac{1}{|\mathcal G|}
 \sum_{g\in\mathcal G}
 \mathds{1}\{T(gx)\ge T(x)\}.
\end{equation*}
对群元素均匀抽样时，每个不同轨道点恰被同样多的稳定子群元素产生，故也等价于对不同轨道点均匀抽样；统计量相同的不同点仍须分别计数。由\eqnrefs{eq:math-stat-group-invariance-null}可知，在原假设下，$X$ 在自己的群轨道上没有一个位置比其他位置更特殊；因此
\begin{equation}\label{eq:math-stat-randomization-validity}
 \mathbb P_{H_0}\{p_{\mathcal G}(X)\le\alpha\}
 \le\alpha.
\end{equation}
离散轨道会使不等号通常严格成立；若在临界秩处再作适当随机化，可以把 水平 调到恰好 $\alpha$。这是有限样本结论，不依赖 中心极限定理、Wilks 或大样本正态近似。
\end{theorem}

\begin{proof}
令 $G$ 在有限群 $\mathcal G$ 上均匀分布，并与 $X$ 独立。由原假设不变性，
\[
 GX\overset{d}=X.
\]
固定样本点 $x$，记群元素产生的统计量多重集为
\[
 \mathcal T_x=\{T(gx):g\in\mathcal G\}.
\]
这里必须保留\emph{多重集}而不是只保留不同轨道点：若 $x$ 有非平凡稳定子群，同一个轨道点会由相同数目的群元素产生。均匀抽取群元素 $G$ 因而恰好等价于从 $\mathcal T_x$ 的 $|\mathcal G|$ 个带重数元素中均匀抽取一个。

令这些数值按非增序排列为
\[
 t_{(1)}\ge t_{(2)}\ge\cdots\ge t_{(M)},
 \qquad M=|\mathcal G|.
\]
对任意群元素 $g$，
\[
 p_{\mathcal G}(gx)
 =\frac1M\#\{h\in\mathcal G:T(hgx)\ge T(gx)\}.
\]
右乘 $g$ 只是群元素的一个置换，所以集合
\[
 \{T(hgx):h\in\mathcal G\}
\]
与 $\mathcal T_x$ 是同一个多重集。若 $T(gx)=t_{(j)}$，则使用“$\ge$”定义尾部时至少有前 $j$ 个元素不小于它，因此
\[
 p_{\mathcal G}(gx)\ge\frac{j}{M}.
\]
于是满足 $p_{\mathcal G}(gx)\le\alpha$ 的群元素至多对应
\[
 j\le \lfloor \alpha M\rfloor
\]
的前 $\lfloor\alpha M\rfloor$ 个位置。即使存在 ties，这个计数只会进一步减少，而不会增加。因此
\[
 \mathbb P\!\left\{
 p_{\mathcal G}(Gx)\le\alpha
 \middle|X=x
 \right\}
 \le
 \frac{\lfloor\alpha M\rfloor}{M}
 \le\alpha.
\]
最后利用 $GX\overset d=X$，
\[
 \mathbb P_{H_0}\{p_{\mathcal G}(X)\le\alpha\}
 =
 \mathbb P_{H_0}\{p_{\mathcal G}(GX)\le\alpha\}
 \le\alpha,
\]
得到\eqnrefs{eq:math-stat-randomization-validity}。因此精确性并不要求群作用自由，也不要求轨道上的统计量没有 ties；真正使用的只有原假设下的群不变性和对群元素的均匀平均。
\end{proof}


\begin{workedexample}[两样本置换与交换性]
最常见的特殊情形是两样本置换检验。设
\[
 X_1,\ldots,X_m,
 \qquad
 Y_1,\ldots,Y_n
\]
在原假设下来自同一分布，并且合并后的 $m+n$ 个观测可交换。条件于合并样本的数值以后，哪 $m$ 个位置被标记为 $X$ 组是等可能的；可取合并位置的对称群 $S_{m+n}$；它作用在容量为 $m$ 的标签子集上，得到 $\binom{m+n}{m}$ 种等可能分组。分组集合本身不是一个群。于是均值差、秩和或其他预先指定统计量的置换分布给出精确条件原假设分布。这里真正的假设是\emph{exchangeability}；若两组在原假设下方差、相关结构或抽样机制仍不同，机械交换标签并不产生精确检验。
\end{workedexample}

\begin{workedexample}[符号检验与符号秩检验]
配对资料给出另一个清楚的群作用。令差值 $D_i=X_i-Y_i$。若 $D_1,\ldots,D_n$ 独立、各自关于 $0$ 对称，并且 $\mathbb P(D_i=0)=0$，则条件于绝对值 $|D_i|$，符号向量在 $\{\pm1\}^n$ 上均匀，故所有 $2^n$ 个独立符号翻转构成精确随机化群。若 $|D_i|$ 出现 ties，可以在给定绝对值后固定使用 midrank 或其他预先规定的分数；精确性来自符号翻转不变性，而不是“秩恰好互异”。以
\[
 T_{\rm SR}
 =\sum_{i=1}^n R_i\mathds{1}\{D_i>0\},
\]
其中 $R_i$ 是 $|D_i|$ 的秩，就得到 Wilcoxon signed-秩 的精确原假设分布。若只假设 $\mathbb P(D_i>0)=1/2$ 且无零差值，则正号个数本身仍服从 $\mathrm{Bin}(n,1/2)$，但 signed-秩 所需的符号--绝对值不变性不再自动成立。这两个检验的假设层级不能混为一谈。
\end{workedexample}

实验设计中的 randomization test 还要再区分一层：若处理标签是由已知随机分配机制产生的，精确性来自\emph{设计本身}的随机化分布；观测值交换性的模型假设并非必要。置换检验与随机化检验可以得到相同计算式，但它们的概率来源不同，因此失效条件也不同。

\section{LAN 下的 LR、Wald 与 Score：同一个法向距离}

\begin{definition}[局部二次似然实验]
设 $\theta\in\RR^d$，真值为 $\theta_0$，模型满足 \chapref{chap:math-stat-estimation-likelihood} 的 局部渐近正态。为看清最基本结构，一维形式为
\begin{equation*}\ell_n\!\left(\theta_0+\frac h{\sqrt n}\right)-\ell_n(\theta_0)
 =h\Delta_n-\frac12h^2I_0+o_{P_{\theta_0}}(1),
 \qquad
 \Delta_n\xrightarrow{d}\mathcal N(0,I_0).
\end{equation*}
局部 似然 是开口向下的随机抛物线；它的顶点同时决定 极大似然估计 的一阶位置与所有二次型检验的一阶统计量。
\end{definition}

\begin{workedexample}[直接验证 Bernoulli 的局部渐近正态性]
在 Bernoulli 模型内点 \(p_0\in(0,1)\)，令
\(p=p_0+h/\sqrt n\)、\(S_n=\sum X_i\)。对
\(\ell_n(p)=S_n\log p+(n-S_n)\log(1-p)\) 作 Taylor 展开，得到
\[
\Delta_n=\frac{S_n-np_0}{\sqrt n\,p_0(1-p_0)},\qquad
I(p_0)=\frac1{p_0(1-p_0)}.
\]
确切地，线性项为 \(h\Delta_n\)。二次导数除以 \(n\) 依大数律趋于
\(-I(p_0)\)，所以二次项趋于 \(-h^2I(p_0)/2\)。
在 \(p_0\) 的一个闭小邻域里三阶导数的绝对值由常数乘 \(n\) 控制；
Taylor 余项因此为
\(O(n|h/\sqrt n|^3)=O(n^{-1/2})\)，对有界 \(h\) 一致趋零。
最后，由中心极限定理 \(\Delta_n\Rightarrow N(0,I(p_0))\)。
这逐步验证了定义，而不是把形式上的二阶展开直接叫作定理。
\end{workedexample}

\begin{proposition}[Fisher 法向投影]
现在考虑光滑复合原假设
\[
 H_0:g(\theta)=0,
 \qquad
 g:\RR^d\to\RR^r,
\]
并设
\[
 G_0=Dg(\theta_0)
\]
满行秩 $r$。零假设在 $\theta_0$ 的切空间为
\[
 \mathcal T_0
 =\{h\in\RR^d:G_0h=0\}.
\]
Fisher 信息定义局部内积
\[
 \langle a,b\rangle_{I_0}
 =a^{\mathsf T}I_0b.
\]
局部渐近正态 的不受限局部极大点为
\[
 h^*=I_0^{-1}\Delta_n.
\]
受限局部极大点则是 $h^*$ 到 $\mathcal T_0$ 的 $I_0$-正交投影。于是原假设可被数据区分的那一部分，恰好是 $h^*$ 在 Fisher 法向空间上的分量。

这个投影可以显式写出。对任意 $h$，其法向分量为
\begin{equation}\label{eq:math-stat-fisher-normal-projection}
 h_\perp
 =I_0^{-1}G_0^{\mathsf T}
 (G_0I_0^{-1}G_0^{\mathsf T})^{-1}
 G_0h.
\end{equation}
它满足 $h-h_\perp\in\mathcal T_0$，并且对任意 $v\in\mathcal T_0$，
\[
 \langle h_\perp,v\rangle_{I_0}=0.
\]
相应平方长度为
\begin{equation}\label{eq:math-stat-fisher-normal-distance}
 \|h_\perp\|_{I_0}^2
 =(G_0h)^{\mathsf T}
 (G_0I_0^{-1}G_0^{\mathsf T})^{-1}
 (G_0h).
\end{equation}
\end{proposition}

\begin{theorem}[正则 LR、Wald 与得分的同一极限]

\label{thm:math-testing-lan-equivalence}
设 $I_0\succ0$，$g$ 在 $\theta_0$ 附近二阶连续可微且 $G_0$ 满行秩，$g(\theta_0)=0$。
除局部一致 LAN 外，要求非受限和受限极大似然估计存在并为根 $n$ 紧，且
\[
 n^{-1/2}U_n(\theta_0)=\Delta_n+o_P(1),\qquad
 \sup_{\|h\|\le M}\left\| -\frac1n\ell_n''(\theta_0+h/\sqrt n)-I_0\right\|\xrightarrow{P}0
\]
对每个 $M<\infty$ 成立，$I(\theta)$ 在真值连续。
这些导数条件保证下述得分统计量与 LAN 中心序列相连；仅有函数值的 LAN 展开不足以直接替换导数。
设 $\widehat\theta_n$ 为不受限 极大似然估计，$\widetilde\theta_n$ 为 $\Theta_0$ 上的受限 极大似然估计。似然-比 统计量定义为
\[
 \operatorname{LR}_n
 =2\{\ell_n(\widehat\theta_n)-\ell_n(\widetilde\theta_n)\}.
\]
Wald 统计量写成
\[
 W_n
 =n\,g(\widehat\theta_n)^{\mathsf T}
 \left[
 G(\widehat\theta_n)
 I(\widehat\theta_n)^{-1}
 G(\widehat\theta_n)^{\mathsf T}
 \right]^{-1}
 g(\widehat\theta_n).
\]
得分 统计量应当只测量受限点处 得分 的法向部分；在光滑约束下可以写为与有效 得分 等价的二次型。若用全参数 得分 表示，则受限极大化的一阶条件已经使切向分量消失，因此
\[
 \mathcal S_n
 =\frac1nU_n(\widetilde\theta_n)^{\mathsf T}
 I(\widetilde\theta_n)^{-1}
 U_n(\widetilde\theta_n)
\]
在一阶上就是同一法向长度。

在上述假设下，三者满足
\begin{equation}\label{eq:math-stat-lr-wald-score-equivalence}
 \operatorname{LR}_n=W_n+o_{\mathbb P}(1)
 =\mathcal S_n+o_{\mathbb P}(1)
 \xrightarrow{d}\chi^2_r.
\end{equation}

自由度 $r$ 是零假设在参数空间中的余维。
\end{theorem}
\begin{proof}
令 $h^*=I_0^{-1}\Delta_n$、$\Pi_\perp$ 为前述 Fisher 法向投影。
局部一致余项、两个极大点的根 $n$ 紧性及二次型的严格凹性给
\[
 \sqrt n(\widehat\theta_n-\theta_0)=h^*+o_P(1),\quad
 \sqrt n(\widetilde\theta_n-\theta_0)=(I-\Pi_\perp)h^*+o_P(1).
\]
配方后，极大对数似然差的两倍为 $\|\Pi_\perp h^*\|_{I_0}^2+o_P(1)$。
对 $g(\widehat\theta_n)$ 的一阶展开也给相同的 Wald 二次型。
局部 Hessian 条件和受限一阶条件则给
\[
 n^{-1/2}U_n(\widetilde\theta_n)
 =\Delta_n-I_0\sqrt n(\widetilde\theta_n-\theta_0)+o_P(1)
 =I_0\Pi_\perp h^*+o_P(1).
\]
因而得分二次型也相同。$I_0^{-1/2}\Delta_n$ 趋于标准 $d$ 维正态，
其在秩 $r$ 的欧氏法向投影上的平方长度服从 $\chi^2_r$。

LR、Wald 和得分的主项可以在同一法向子空间中分别计算。

局部渐近正态 主项为
\[
 q_n(h)
 =h^{\mathsf T}\Delta_n-\frac12h^{\mathsf T}I_0h
 =-\frac12\|h-h^*\|_{I_0}^2
 +\frac12\|h^*\|_{I_0}^2,
\]
其中 $h^*=I_0^{-1}\Delta_n$。受限极大点是
\[
 \widetilde h^*=h^*-h_\perp^*,
\]
其中由\eqnrefs{eq:math-stat-fisher-normal-projection}
\[
 h_\perp^*
 =I_0^{-1}G_0^{\mathsf T}
 (G_0I_0^{-1}G_0^{\mathsf T})^{-1}
 G_0h^*.
\]
所以 LR 主项为
\[
 2\{q_n(h^*)-q_n(\widetilde h^*)\}
 =\|h_\perp^*\|_{I_0}^2.
\]
另一方面，由约束的一阶展开和 极大似然估计 渐近线性表示，
\[
 \sqrt n\,g(\widehat\theta_n)
 =G_0h^*+o_{\mathbb P}(1),
\]
代入 Wald 统计量并使用\eqnrefs{eq:math-stat-fisher-normal-distance}，得到
\[
 W_n
 =\|h_\perp^*\|_{I_0}^2+o_{\mathbb P}(1).
\]
对 得分，在受限 极大似然估计 处考虑 Lagrange 条件
\[
 U_n(\widetilde\theta_n)
 =G(\widetilde\theta_n)^{\mathsf T}\kappa_n.
\]
这说明 得分 落在约束法向余切空间。将 $U_n/\sqrt n$ 用 局部渐近正态 线性化并乘以 $I_0^{-1}$，得到的参数空间向量正是 $h_\perp^*+o_{\mathbb P}(1)$，故 得分 也给出相同平方长度。最后
\[
 I_0^{-1/2}\Delta_n\xrightarrow{d}\mathcal N_d(0,I_d),
\]
而标准高斯向量在 $r$ 维法向子空间上的平方长度服从 $\chi^2_r$，于是得到\eqnrefs{eq:math-stat-lr-wald-score-equivalence}。

在已知方差的正态均值检验 $H_0:\mu=0$ 对 $H_1:\mu\neq0$ 中，LR、Wald 与 得分 统计量都等于 $n\bar X_n^2/\sigma^2$，所以三者精确相同。一般模型只保证它们在正则条件下渐近等价；有限样本中，Wald 统计量还会随非线性重参数化改变。
\end{proof}


\begin{exbox}[Bernoulli 模型中 LR、Wald 与 得分 的有限样本差异]
令 $X_i\iid\mathrm{Bernoulli}(p)$，检验 $H_0:p=p_0$，其中 $0<p_0<1$，并记 $\widehat p=\overline X$。当 $0<\widehat p<1$ 时，三类统计量分别为
\begin{align*}
 W_n
 &=\frac{n(\widehat p-p_0)^2}
 {\widehat p(1-\widehat p)},\\
 \mathcal S_n
 &=\frac{n(\widehat p-p_0)^2}
 {p_0(1-p_0)},\\
 \operatorname{LR}_n
 &=2n\left[
 \widehat p\log\frac{\widehat p}{p_0}
 +(1-\widehat p)\log\frac{1-\widehat p}{1-p_0}
 \right].
\end{align*}
它们在有限样本中不是同一个函数：Wald 在 极大似然估计 处估计曲率，得分 在原假设处估计曲率，LR 则直接比较两个极大 似然。若 $H_0$ 成立，$\widehat p-p_0=O_{\mathbb P}(n^{-1/2})$。对 Bernoulli KL 距离在 $p_0$ 处作 Taylor 展开，
\[
 2\left[
 p\log\frac p{p_0}
 +(1-p)\log\frac{1-p}{1-p_0}
 \right]
 =\frac{(p-p_0)^2}{p_0(1-p_0)}
 +O(|p-p_0|^3).
\]
因此
\[
 \operatorname{LR}_n-\mathcal S_n=O_{\mathbb P}(n^{-1/2}),
 \qquad
 W_n-\mathcal S_n=o_{\mathbb P}(1),
\]
三者才在一阶上共同趋于 $\chi^2_1$。若样本全为 $0$ 或全为 $1$，则 $\widehat p$ 落在边界，以上 Wald 形式甚至因 $\widehat p(1-\widehat p)=0$ 而失效，而 得分 与 LR 在 $0<p_0<1$ 下仍有良好定义。这进一步说明 Wald 统计量对参数坐标和边界更敏感；“三者正则渐近等价”绝不是有限样本可互换性。
\end{exbox}
\begin{proposition}[局部备择的非中心卡方与功效]
沿用上一检验定理的局部正则条件，并要求所有 $o_P(1)$ 余项在有界局部参数集上一致。
对固定备择 $\theta\ne\theta_0$，一致检验的功效往往趋于 $1$，因此不能描述统计程序在极限中的最细分辨率。局部渐近正态 的自然局部备择是
\[
 \theta_{n,h}
 =\theta_0+\frac h{\sqrt n}.
\]
在 $P_{\theta_0}$ 下，对数似然比 满足
\[
 \log\frac{\dd P_{\theta_{n,h}}^{(n)}}
 {\dd P_{\theta_0}^{(n)}}
 =h^{\mathsf T}\Delta_n
 -\frac12h^{\mathsf T}I_0h
 +o_{\mathbb P}(1).
\]
其极限是均值为 $-h^{\mathsf T}I_0h/2$、方差为 $h^{\mathsf T}I_0h$ 的正态变量。这里若一列事件 $A_n$ 满足 $P_{\theta_0}^{(n)}(A_n)\to0$ 时也必有 $P_{\theta_{n,h}}^{(n)}(A_n)\to0$，就称后者关于前者 contiguous；若两个方向都成立，则称两列实验 相互 contiguous。局部渐近正态 的极限似然比严格为正且期望为 $1$，由 Le Cam 第一引理得到这里的 mutual contiguity。它保证 $n^{-1/2}$ 局部换测度不会把原假设下的零概率极限事件突然变成正概率事件。

Le Cam 第三引理把测度变化作用到中心序列上，得到在 $P_{\theta_{n,h}}$ 下
\begin{equation}\label{eq:math-stat-third-lemma-shift}
 \Delta_n
 \xrightarrow{d}\mathcal N_d(I_0h,I_0).
\end{equation}
所以标准化向量 $I_0^{-1/2}\Delta_n$ 的均值从 $0$ 平移到 $I_0^{1/2}h$。把它投影到零假设法向空间后，LR/Wald/得分 极限都变成
\[
 \chi^2_r(\Lambda_h),
\]
其中
\begin{equation*}\Lambda_h
 =\|h_\perp\|_{I_0}^2
 =(G_0h)^{\mathsf T}
 (G_0I_0^{-1}G_0^{\mathsf T})^{-1}
 (G_0h).
\end{equation*}
因此水平 $\alpha$ 检验的局部渐近功效为
\[
 \mathbb P\{\chi^2_r(\Lambda_h)>\chi^2_{r,1-\alpha}\}.
\]
沿零假设切空间的 $h$ 有 $\Lambda_h=0$，一阶上不可区分；只有 Fisher 法向分量贡献局部功效。
\end{proposition}

\begin{proof}
在原假设下，联合向量
\[
 \left(
 \Delta_n,
 \log\frac{\dd P_{\theta_{n,h}}^{(n)}}
 {\dd P_{\theta_0}^{(n)}}
 \right)
\]
由 局部渐近正态 渐近等价于
\[
 \left(
 Z,
 h^{\mathsf T}Z-\frac12h^{\mathsf T}I_0h
 \right),
 \qquad
 Z\sim\mathcal N_d(0,I_0).
\]
换到局部备择测度 $P_{\theta_{n,h}}^{(n)}$，相当于用指数权重
\[
 \exp\!\left(
 h^{\mathsf T}Z-\frac12h^{\mathsf T}I_0h
 \right)
\]
重新加权原高斯分布。
为直接写成密度比，先取局部模型互相绝对连续的共同正支撑情形。
LAN 似然比的非负极限具有期望 $1$，加上原似然比期望 $1$，给出一致可积性。
因此对任意有界连续 $\varphi$，联合弱收敛与一致可积性给

\[
\mathbb E_{\theta_{n,h}}\varphi(\Delta_n)
=\mathbb E_{\theta_0}\!\left[
\varphi(\Delta_n)\frac{\dd P_{\theta_{n,h}}^{(n)}}{\dd P_{\theta_0}^{(n)}}
\right]
\longrightarrow\mathbb E\!\left\{
\varphi(Z)\exp\!\left(h^{\mathsf T}Z-\frac12h^{\mathsf T}I_0h\right)
\right\}.
\]
把 $Z$ 的密度 $(2\pi)^{-d/2}|I_0|^{-1/2}\exp(-\frac12z^{\mathsf T}I_0^{-1}z)$ 与权重相乘，指数合并为
\[
 -\frac12z^{\mathsf T}I_0^{-1}z
 +h^{\mathsf T}z
 -\frac12h^{\mathsf T}I_0h.
\]
完成平方：令 $w=z-I_0h$，则
\[
 -\frac12z^{\mathsf T}I_0^{-1}z
 +h^{\mathsf T}z
 -\frac12h^{\mathsf T}I_0h
 =-\frac12(z-I_0h)^{\mathsf T}
 I_0^{-1}(z-I_0h),
\]
因此加权后的 $Z$ 密度正比于 $\exp(-\frac12w^{\mathsf T}I_0^{-1}w)$，分布正是 $\mathcal N_d(I_0h,I_0)$，这就得到\eqnrefs{eq:math-stat-third-lemma-shift}。
\end{proof}


\begin{exbox}[正态均值设计把功效写成标准化法向距离]
设 $X_i\iid\mathcal N(\mu,\sigma^2)$，其中 $\sigma$ 已知，检验
\[
 H_0:\mu\le\mu_0,
 \qquad
 H_1:\mu>\mu_0.
\]
由单调 似然比，水平 $\alpha$ 的 一致最强检验在
\[
 Z_{\mu,n}=\frac{\sqrt n(\overline X-\mu_0)}{\sigma}
 >z_{1-\alpha}
\]
时拒绝。对任意固定 $\mu_1>\mu_0$，其有限样本精确功效为
\[
 \operatorname{Power}_n(\mu_1)
 =1-\Phi\!\left(
 z_{1-\alpha}
 -\frac{\sqrt n(\mu_1-\mu_0)}{\sigma}
 \right).
\]
若希望该点的功效至少为 $1-\beta$，等价于
\[
 \frac{\sqrt n(\mu_1-\mu_0)}{\sigma}
 \ge z_{1-\alpha}+z_{1-\beta},
\]
因而
\[
 n\ge
 \left[
 \frac{\sigma(z_{1-\alpha}+z_{1-\beta})}
 {\mu_1-\mu_0}
 \right]^2.
\]
这不是额外的“样本量公式”：分子中的高斯分位数来自错误概率约束，分母中的效应量通过 $\sqrt n(\mu_1-\mu_0)/\sigma$ 进入，而这个量正是 Fisher 度量下的法向距离。局部渐近正态 中的非中心参数只是把这套精确 Gaussian 几何推广到一般正则模型。
\end{exbox}

\begin{proposition}[干扰参数与有效信息]
将参数分块为
\[
 \theta=(\vartheta,\eta),
\]
其中 $\vartheta\in\RR^r$ 是关注参数，$\eta\in\RR^q$ 是 干扰参数。Fisher 信息写成
\[
 I=
 \begin{pmatrix}
 I_{\vartheta\vartheta}&I_{\vartheta\eta}\\
 I_{\eta\vartheta}&I_{\eta\eta}
 \end{pmatrix}.
\]
若 $I_{\eta\eta}$ 可逆，关注参数 得分 中可由 干扰得分 线性预测的部分应当先投影掉。定义有效 得分
\begin{equation*}s_{\vartheta\cdot\eta}
 =s_\vartheta-I_{\vartheta\eta}I_{\eta\eta}^{-1}s_\eta.
\end{equation*}
由于
\[
 \Cov(s_\vartheta,s_\eta)=I_{\vartheta\eta},
 \qquad
 \Cov(s_\eta)=I_{\eta\eta},
\]
它正是 $s_\vartheta$ 对 干扰得分 线性子空间做 $L^2$ 投影后的残差。其协方差为 Schur 补
\begin{equation*}I_{\vartheta\cdot\eta}
 =I_{\vartheta\vartheta}
 -I_{\vartheta\eta}I_{\eta\eta}^{-1}I_{\eta\vartheta}.
\end{equation*}
剖面似然、受限 得分 与 Wald 协方差在一阶上都使用同一个有效信息。若 $I_{\vartheta\cdot\eta}$ 奇异，说明某些关注方向在消去 冗余 后已没有一阶可识别信息，普通 $\sqrt n$ 正则推断必须失效。
\end{proposition}

\begin{proof}
把 局部渐近正态 局部参数和中心序列按同样方式分块：
\[
 h=\begin{pmatrix}h_\vartheta\\h_\eta\end{pmatrix},
 \qquad
 \Delta_n=\begin{pmatrix}\Delta_{\vartheta,n}\\\Delta_{\eta,n}\end{pmatrix}.
\]
局部二次型写成
\begin{align*}
 q_n(h_\vartheta,h_\eta)
 &=h_\vartheta^{\mathsf T}\Delta_{\vartheta,n}
   +h_\eta^{\mathsf T}\Delta_{\eta,n}\\
 &\quad-\frac12h_\vartheta^{\mathsf T}I_{\vartheta\vartheta}h_\vartheta
   -h_\vartheta^{\mathsf T}I_{\vartheta\eta}h_\eta
   -\frac12h_\eta^{\mathsf T}I_{\eta\eta}h_\eta.
\end{align*}
对固定 $h_\vartheta$，关于 $h_\eta$ 的一阶条件给出
\[
 \widehat h_\eta(h_\vartheta)
 =I_{\eta\eta}^{-1}
 \{\Delta_{\eta,n}-I_{\eta\vartheta}h_\vartheta\}.
\]
把它代回 $q_n$，与 $h_\vartheta$ 无关的项并入常数 $C_n$，得到 剖面局部渐近正态：
\begin{align*}
 q_n^{\rm prof}(h_\vartheta)
 &=C_n
 +h_\vartheta^{\mathsf T}
 \underbrace{\left(
 \Delta_{\vartheta,n}
 -I_{\vartheta\eta}I_{\eta\eta}^{-1}\Delta_{\eta,n}
 \right)}_{\Delta_{\vartheta\cdot\eta,n}}\\
 &\quad-\frac12h_\vartheta^{\mathsf T}
 \underbrace{\left(
 I_{\vartheta\vartheta}
 -I_{\vartheta\eta}I_{\eta\eta}^{-1}I_{\eta\vartheta}
 \right)}_{I_{\vartheta\cdot\eta}}
 h_\vartheta.
\end{align*}
因此 干扰参数并不是“把协方差矩阵某一块删掉”。必须先从关注 得分 中减去它在 干扰得分 线性空间上的 $L^2$ 投影：
\[
 \Delta_{\vartheta\cdot\eta,n}
 =\Delta_{\vartheta,n}
 -I_{\vartheta\eta}I_{\eta\eta}^{-1}\Delta_{\eta,n}.
\]
直接计算协方差，
\[
 \Cov(\Delta_{\vartheta\cdot\eta,n})
 \to I_{\vartheta\cdot\eta},
\]
所以 Schur 补 正是投影后剩余的 Fisher 信息。剖面 LR、有效 得分 与关注参数的 Wald 协方差之所以一阶一致，根源就在这个同一个投影后的局部高斯实验。
\end{proof}


\begin{note}[边界、参数依赖支撑与奇异模型]
\eqnrefs{eq:math-stat-lr-wald-score-equivalence} 依赖真值位于光滑参数空间内部、局部可识别、Fisher 信息非奇异以及 二次均值可微/局部渐近正态。若这些条件失效，$\chi^2$ 极限没有理由保持不变。

最简单的边界例子是单参数 $\theta\ge0$ 且检验 $\theta=0$。局部无约束高斯极大点若为负，会被边界截到 $0$，典型 LR 极限成为
\[
 \frac12\delta_0+\frac12\chi^2_1,
\]
而不是 $\chi^2_1$。若 $X_i\iid\mathrm{Unif}(0,\theta)$，最大次序统计量以 $n^{-1}$ 而非 $n^{-1/2}$ 速度逼近端点，支撑依赖参数使普通 得分/Fisher 展开失效。有限混合模型在成分重合处还会同时出现不可识别与奇异信息矩阵。遇到这些模型时，应先重新确定局部尺度和极限实验，而不是机械查 Wilks 自由度。
\end{note}

\section{Pearson 卡方：概率单纯形上的切空间投影}

\begin{definition}[单纯形中的标准化计数波动]
设
\[
 (N_1,\ldots,N_k)
 \sim\operatorname{Multinomial}(n;p_1,\ldots,p_k),
 \qquad
 p_j>0,
 \qquad
 \sum_{j=1}^kp_j=1.
\]
记
\[
 p=(p_1,\ldots,p_k)^{\mathsf T},
 \qquad
 D_p=\diag(p_1,\ldots,p_k),
 \qquad
 s=(\sqrt{p_1},\ldots,\sqrt{p_k})^{\mathsf T}.
\]
定义标准化频数残差
\begin{equation*}Z_n
 =D_p^{-1/2}\frac{N-np}{\sqrt n}.
\end{equation*}
因为 $\sum_j(N_j-np_j)=0$，
\[
 s^{\mathsf T}Z_n=0,
\]
所以 $Z_n$ 始终位于 $s^\perp$，这正是概率单纯形在 $p$ 处经过方差标准化后的 $(k-1)$ 维切空间。
\end{definition}


\begin{theorem}[固定正概率的 Pearson 极限]
设 $k\ge2$ 固定、$p_j>0$ 固定，$N\sim\operatorname{Multinomial}(n,p)$，则
\[
 \sum_j\frac{(N_j-np_j)^2}{np_j}\xrightarrow{d}\chi^2_{k-1}.
\]
\end{theorem}
\begin{proof}
把每次分类观测写成 one-hot 向量 $Y_i\in\{e_1,\ldots,e_k\}$，则
\[
 N=\sum_{i=1}^nY_i,
 \qquad
 \mathbb E Y_i=p,
 \qquad
 \Cov(Y_i)=D_p-pp^{\mathsf T}.
\]
于是
\begin{equation*}\Cov(Z_n)
 =D_p^{-1/2}(D_p-pp^{\mathsf T})D_p^{-1/2}
 =I_k-ss^{\mathsf T}.
\end{equation*}
右端恰好是投影到 $s^\perp$ 的正交投影。多元 中心极限定理 因此给出退化高斯极限，它在该切空间上是标准球对称正态。

Pearson 统计量就是这个向量的平方长度：
\begin{equation}\label{eq:math-stat-pearson-quadratic-form}
 X_n^2
 =\sum_{j=1}^k\frac{(N_j-np_j)^2}{np_j}
 =Z_n^{\mathsf T}Z_n.
\end{equation}
所以当 $p_j$ 固定为正时
\begin{equation*}X_n^2\xrightarrow{d}\chi^2_{k-1}.
\end{equation*}
自由度是切空间维数。对任意固定 $t\in\RR^k$，$t^{\mathsf T}(Y_i-p)$ 是独立同分布有界变量；标量 CLT 与 Cramér--Wold 准则给出上述向量高斯极限，连续映射再给平方长度极限。
\end{proof}

\begin{theorem}[拟合概率的 Pearson 极限]
设 $p(\eta)$ 在真值附近二阶连续可微，$p_j(\eta_0)>0$；使用一致的内点多项分布 MLE，标准化 Jacobian 满列秩 $q<k-1$。
更重要的是原假设本身含待估参数的情形。设
\[
 p=p(\eta),
 \qquad
 \eta\in\RR^q,
\]
且 $\eta_0$ 是真值。记 $p_0=p(\eta_0)$、$D_{p_0}=\diag(p_0)$，并令
\[
 \dot p_0
 =\left.
 \frac{\partial p(\eta)}{\partial\eta^{\mathsf T}}
 \right|_{\eta_0}
 \in\RR^{k\times q}.
\]
由于 $\mathbf1^{\mathsf T}p(\eta)=1$，有
\[
 \mathbf1^{\mathsf T}\dot p_0=0.
\]
若 $\dot p_0$ 满列秩 $q<k-1$，标准化模型切空间由
\[
 B_{p,0}=D_{p_0}^{-1/2}\dot p_0
\]
的列空间张成，并且 $s^{\mathsf T}B_{p,0}=0$，所以 $\mathcal C(B_{p,0})$ 是 $s^\perp$ 内的 $q$ 维子空间。拟合 $\eta$ 会吸收经验频率涨落在这些模型切向方向上的分量。

此时真正用于拟合优度检验的 Pearson 统计量必须以拟合概率为中心并用拟合概率标准化：
\[
 X_{n,\mathrm{fit}}^2
 =\sum_{j=1}^k
 \frac{\{N_j-np_j(\widehat\eta)\}^2}
 {np_j(\widehat\eta)}.
\]
它与前面固定 $p$ 的 $X_n^2$ 不是同一个随机变量；“估计参数后减自由度”也不能直接套到固定中心的\eqnrefs{eq:math-stat-pearson-quadratic-form}上。在上述内部点、满秩与光滑正则条件下，得分 方程恰好消去 $s^\perp$ 中的 $q$ 个 Fisher 切向方向，因此剩余涨落位于维数 $k-1-q$ 的正交补中，并有
\begin{equation}\label{eq:math-stat-pearson-fitted-df}
 X_{n,\mathrm{fit}}^2\xrightarrow{d}\chi^2_{k-1-q}.
\end{equation}
\end{theorem}


\begin{proof}
在此证明内记 $D=\operatorname{diag}(p_j(\eta_0))$、$J=\dot p_0$、$A=D^{-1/2}J$、
$v=(\sqrt{p_j})_j$、$r=q$，并用 $\theta=\eta$、$\theta_0=\eta_0$ 表示参数。
假设使用内点多项分布极大似然估计 \(\widehat\theta\)，且它一致。对数似然得分在真值的第 \(a\) 分量为
\[
U_a(\theta_0)
=\sum_i\frac{N_i}{p_i}J_{ia}
=\sum_i\frac{N_i-np_i}{p_i}J_{ia}
=\sqrt n\,(A^\top Z_n)_a;
\]
第二个等号正用到 \(\sum_iJ_{ia}=0\)。其负 Hessian 除以 \(n\) 依概率趋于
\(J^\top D^{-1}J=A^\top A\)。在 \(\widehat\theta\) 处令得分为零并作 Taylor 展开，便得到
\[
\sqrt n(\widehat\theta-\theta_0)
=(A^\top A)^{-1}A^\top Z_n+o_P(1).
\]
将 \(p(\widehat\theta)=p+J(\widehat\theta-\theta_0)+o_P(n^{-1/2})\)
代入拟合后的标准化残差，得到
\[
D^{-1/2}\sqrt n\{N/n-p(\widehat\theta)\}
=\{I-A(A^\top A)^{-1}A^\top\}Z_n+o_P(1).
\]
括号中减去的是投到 \(\operatorname{col}(A)\) 的正交投影；由于这 \(r\) 个方向位于 \(v^\perp\)，剩余投影在 \(v^\perp\) 上的秩是 \(k-1-r\)。把分母中的 \(p_i(\widehat\theta)\) 以一致性换成 \(p_i\)，差为 \(o_P(1)\)，故拟合后的 Pearson 统计量趋于
\(\chi^2_{k-1-r}\)。边界、零概率或秩亏须重新分析；饱和情形 \(r=k-1\) 没有剩余检验方向，形式极限为退化的 \(\chi^2_0\)。

在多项极大似然估计下，得分方程将切空间投影逐项写为以下形式。

记经验概率 $\widehat p=N/n$，并令 $p_0=p(\eta_0)$、$D_{p_0}=\diag(p_0)$。这里取 $\widehat\eta$ 为位于参数空间内部的 多项 极大似然估计，并假设 $p_j(\eta_0)>0$、$\dot p_0$ 满列秩且 $p(\eta)$ 在 $\eta_0$ 附近二阶可微。模型的 对数似然 为
\[
 \ell_n(\eta)=\sum_{j=1}^kN_j\log p_j(\eta).
\]
因为 $\mathbf1^{\mathsf T}\dot p(\eta)=0$，得分 可写成
\begin{align*}
 U_n(\eta)
 &=\dot p(\eta)^{\mathsf T}
 D_{p(\eta)}^{-1}N\\
 &=\dot p(\eta)^{\mathsf T}
 D_{p(\eta)}^{-1}
 \{N-np(\eta)\}.
\end{align*}
在真值处令
\[
 Z_n=D_{p_0}^{-1/2}\sqrt n(\widehat p-p_0),
 \qquad
 B_{p,0}=D_{p_0}^{-1/2}\dot p_0.
\]
于是
\[
 \frac1{\sqrt n}U_n(\eta_0)
 =B_{p,0}^{\mathsf T}Z_n.
\]
另一方面，单位样本 Fisher 信息为
\[
 I_\eta(\eta_0)
 =\dot p_0^{\mathsf T}D_{p_0}^{-1}\dot p_0
 =B_{p,0}^{\mathsf T}B_{p,0}.
\]
对 得分 方程 $U_n(\widehat\eta)=0$ 在 $\eta_0$ 处作 Taylor 展开，并使用 $-n^{-1}\nabla U_n(\eta_0)\xrightarrow{\mathbb P} I_\eta(\eta_0)$，得到
\[
 0
 =B_{p,0}^{\mathsf T}Z_n
 -B_{p,0}^{\mathsf T}B_{p,0}
 \sqrt n(\widehat\eta-\eta_0)
 +o_{\mathbb P}(1).
\]
由于 $B_{p,0}$ 满列秩，
\[
 \sqrt n(\widehat\eta-\eta_0)
 =(B_{p,0}^{\mathsf T}B_{p,0})^{-1}B_{p,0}^{\mathsf T}Z_n
 +o_{\mathbb P}(1).
\]
再乘以 $B_{p,0}$，便得到拟合掉的标准化频率分量
\[
 B_{p,0}\sqrt n(\widehat\eta-\eta_0)
 =P_{B_{p,0}}Z_n+o_{\mathbb P}(1),
 \qquad
 P_{B_{p,0}}
 =B_{p,0}(B_{p,0}^{\mathsf T}B_{p,0})^{-1}B_{p,0}^{\mathsf T}.
\]
同时，参数映射的一阶展开给出
\[
 D_{p_0}^{-1/2}\sqrt n
 \{p(\widehat\eta)-p_0\}
 =P_{B_{p,0}}Z_n+o_{\mathbb P}(1).
\]
因此以真值概率标准化的拟合残差为
\[
 D_{p_0}^{-1/2}\sqrt n
 \{\widehat p-p(\widehat\eta)\}
 =(I-P_{B_{p,0}})Z_n+o_{\mathbb P}(1).
\]
由于 $\widehat\eta\xrightarrow{\mathbb P}\eta_0$ 且所有 $p_j(\eta_0)>0$，把分母中的 $D_{p_0}$ 换成 $D_{p(\widehat\eta)}$ 只改变 $o_{\mathbb P}(1)$ 项，所以真正的 $X_{n,\mathrm{fit}}^2$ 具有同一极限。

原始 多项 涨落满足 $Z_n\xrightarrow{d}\mathcal N_k(0,I-ss^{\mathsf T})$，其中 $s=(\sqrt{p_{01}},\ldots,\sqrt{p_{0k}})^{\mathsf T}$，故极限随机向量位于 $s^\perp$。又因为 $s^{\mathsf T}B_{p,0}=0$，$\mathcal C(B_{p,0})$ 是 $s^\perp$ 内的 $q$ 维子空间。投影掉它之后，剩余球对称 Gaussian 分量位于维数
\[
 (k-1)-q
\]
的正交补中，于是得到\eqnrefs{eq:math-stat-pearson-fitted-df}。所以“估计 $q$ 个参数就扣掉 $q$ 个自由度”不是计数口诀，而是 得分 方程把经验涨落沿模型 Fisher 切空间的 $q$ 个方向正交消去后的几何结果。
\end{proof}

这里的自由度下降是投影维数的结论；若真值位于边界、模型切空间退化或某些格概率趋于零，则这个正则卡方极限不能机械沿用。

\begin{proposition}[多项似然偏差与 Pearson 的一阶等价]
对于固定的完全指定概率向量 $p$，saturated model 与 $p$ 之间的 似然-比 goodness-of-fit 统计量为
\[
 \mathcal D_n^{2}
 =2\sum_{j:N_j>0}
 N_j\log\frac{N_j}{np_j}.
\]
写 $N_j=np_j(1+\delta_j)$，其中 $\delta_j=O_{\mathbb P}(n^{-1/2})$。利用
\[
 (1+\delta)\log(1+\delta)
 =\delta+\frac12\delta^2+O(\delta^3)
\]
以及 $\sum_jnp_j\delta_j=\sum_j(N_j-np_j)=0$，线性项抵消，得到
\[
 \mathcal D_n^{2}
 =\sum_jnp_j\delta_j^2+o_{\mathbb P}(1)
 =X_n^2+o_{\mathbb P}(1).
\]
所以完全指定 $p$ 时 Pearson 与 多项 LR 的一阶等价来自同一个局部二次 似然，而不是偶然共享 $\chi^2$ 表。

若零假设是前面的光滑参数族 $p(\eta)$，LR 也必须把参数拟合计入，而不能继续把真值 $p_0$ 当作已知。定义 fitted 偏差
\begin{equation*}\mathcal D_{n,\mathrm{fit}}^{2}
 =2\sum_{j:N_j>0}
 N_j\log\frac{N_j}{n p_j(\widehat\eta)}.
\end{equation*}
令
\[
 \delta_{j,n}^{\rm fit}
 =\frac{N_j-np_j(\widehat\eta)}{np_j(\widehat\eta)}.
\]
正则 极大似然估计 的 $n^{-1/2}$ 局部化与 $p_j(\eta_0)>0$ 给出 $\max_j|\delta_{j,n}^{\rm fit}|=O_{\mathbb P}(n^{-1/2})$。而
\[
 \sum_jnp_j(\widehat\eta)\delta_{j,n}^{\rm fit}
 =\sum_jN_j-n=0.
\]
对每一项使用同一个 Taylor 展开，有限个 单元 的三阶余项总和为 $o_{\mathbb P}(1)$，因此
\begin{equation*}\mathcal D_{n,\mathrm{fit}}^{2}
 =X_{n,\mathrm{fit}}^2+o_{\mathbb P}(1)
 \xrightarrow{d}\chi^2_{k-1-q}.
\end{equation*}
所以“拟合 $q$ 个参数后”的 Pearson 与 LR 仍由同一个 $k-1-q$ 维法向子空间控制；区别只在有限样本中使用不同的二次近似，而不是自由度规则发生了两套独立来源。
\end{proposition}

\begin{workedexample}[列联表独立性的法向维数]
列联表只是这一结构的特殊参数模型。对 $r\times c$ 表，完整单纯形维数是 $rc-1$；独立性模型
\[
 p_{ij}=a_ib_j
\]
具有
\[
 (r-1)+(c-1)=r+c-2
\]
个自由参数。因此法向维数为
\[
 (rc-1)-(r+c-2)=(r-1)(c-1).
\]
用边际频率的 极大似然估计 得到期望频数
\[
 E_{ij}=\frac{N_{i+}N_{+j}}{n},
\]
于是
\begin{equation*}X^2
 =\sum_{i=1}^r\sum_{j=1}^c
 \frac{(N_{ij}-E_{ij})^2}{E_{ij}}
 \xrightarrow{d}\chi^2_{(r-1)(c-1)}.
\end{equation*}
若存在结构零、约束参数或单元格概率随 $n$ 变小，则模型切空间维数和普通卡方近似都需要重新检查。
\end{workedexample}

\section{正态精确检验与检验--置信集对偶}

局部渐近正态 给出广泛适用的大样本结构，但 \chapref{chap:math-stat-gaussian-sampling} 的正态投影在特定模型中能给更强的有限样本精确检验。两者应该通过模型强弱而不是通过“传统方法/现代方法”区分。

\begin{proposition}[未知方差的 t 枢轴与精确 LR 排序]
取 $n\ge2$，方差严格为正，观测独立正态。
若 $X_i\iid\mathcal N(\mu,\sigma^2)$、$\sigma^2$ 未知，检验
\[
 H_0:\mu=\mu_0
\]
时
\[
 T
 =\frac{\overline X-\mu_0}{S/\sqrt n}
 \sim\mathrm t_{n-1}
\]
在原假设下精确成立。双侧水平 $\alpha$ 检验拒绝条件为
\[
 |T|>t_{n-1,1-\alpha/2}.
\]
它并不是 Wilks 定理的有限样本“近似版”，而是独立 Gaussian 投影 的精确结果。

从 似然比 也能看出二者在这个模型中严格等价。完整模型 极大似然估计 为
\[
 \widehat\mu=\overline X,
 \qquad
 \widehat\sigma^2=\frac1n\sum_i(X_i-\overline X)^2,
\]
原假设下
\[
 \widetilde\sigma^2=\frac1n\sum_i(X_i-\mu_0)^2.
\]
平方和分解给出
\[
 \widetilde\sigma^2
 =\widehat\sigma^2
 \left(1+\frac{T^2}{n-1}\right),
\]
所以 似然比
\[
 \Lambda_{\rm LR}
 =\left(
 1+\frac{T^2}{n-1}
 \right)^{-n/2}
\]
严格随 $|T|$ 递减。因此双侧 LR 与双侧 $t$ 检验具有完全相同的拒绝域排序；这里是有限样本单调变换等价，而不是渐近等价。
\end{proposition}

\begin{proposition}[方差检验的卡方校准]
若均值 $\mu$ 已知并检验
\[
 H_0:\sigma^2=\sigma_0^2,
\]
令
\[
 S_\mu=\sum_i(X_i-\mu)^2,
 \qquad
 \tau=\sigma^2.
\]
对数似然 为
\[
 \ell(\tau)=C-\frac n2\log\tau-\frac{S_\mu}{2\tau},
\]
不受限 极大似然估计 为 $\widehat\tau=S_\mu/n$，从而
\begin{equation*}-2\log\Lambda_{\rm LR}
 =n\log\frac{\sigma_0^2}{S_\mu/n}
 +\frac{S_\mu}{\sigma_0^2}-n.
\end{equation*}
在原假设下 $U=S_\mu/\sigma_0^2\sim\chi^2_n$，右端等于
\[
 U-n-n\log(U/n),
\]
它在 $U=n$ 处最小并向两侧增大。因此 LR 拒绝域同样可以由精确卡方分布校准。
若 $\mu$ 未知，记 $\mathrm{RSS}=\sum_i(X_i-\bar X)^2$，原假设下
$U=\mathrm{RSS}/\sigma_0^2\sim\chi^2_{n-1}$。但剖面 LR 本身仍为
\[
 -2\log\Lambda_{\rm LR}=U-n-n\log(U/n).
\]
自由度 $n-1$ 属于校准 $U$ 的分布，不能把 LR 公式中的所有 $n$ 机械替成 $n-1$。

\end{proposition}

\begin{theorem}[随机化检验与置信集的反演]
最后，检验与置信集具有一般对偶性，但随机化检验必须区分“拒绝概率”与真正发生的拒绝事件。若对每个 $\theta_0$ 都有水平 $\alpha$ 的检验函数 $\phi_{\theta_0}(x)\in[0,1]$，再引入与数据独立的
\[
 U\sim\mathrm{Unif}(0,1),
\]
并规定当 $U\le\phi_{\theta_0}(x)$ 时实际拒绝 $\theta_0$。反演所有未被拒绝的参数值得到随机置信集
\begin{equation}\label{eq:math-stat-randomized-test-confidence-duality}
 C_{1-\alpha}(x,U)
 =\{\theta_0:U>\phi_{\theta_0}(x)\}.
\end{equation}
于是对任意真值 $\theta$，利用 $U$ 的条件均匀性，
\begin{align*}
 \mathbb P_\theta\{\theta\in C_{1-\alpha}(X,U)\}
 &=\mathbb E_\theta\!
 \left[
 \mathbb P\{U>\phi_\theta(X)\mid X\}
 \right]\\
 &=1-\mathbb E_\theta\phi_\theta(X)
 \ge1-\alpha.
\end{align*}
当检验非随机化、$\phi_{\theta_0}\in\{0,1\}$ 时，辅助变量 $U$ 不再影响接受与否，\eqnrefs{eq:math-stat-randomized-test-confidence-duality}才退化为熟悉的
\[
 C_{1-\alpha}(x)=\{\theta_0:\phi_{\theta_0}(x)=0\}.
\]
若检验 水平 恰为 $\alpha$，覆盖率也恰为 $1-\alpha$；若检验只在渐近意义下控制水平，相应置信集也只有渐近覆盖保证。这与\eqnrefs{eq:math-stat-confidence-set-definition}的定义完全一致。
\end{theorem}

\begin{definition}[剖面似然置信集]
对含 干扰参数 的正则 似然，这个对偶性最自然地写成 剖面似然 置信集。把参数分块为 $\theta=(\vartheta,\eta)\in\Theta$，其中 $\vartheta\in\RR^r$ 是关注参数。关注参数的可行空间与给定 $\vartheta$ 时的 干扰参数切片分别记为
\[
 \mathcal V
 =\{\vartheta:\exists\eta,\ (\vartheta,\eta)\in\Theta\},
 \qquad
 \mathcal H(\vartheta)
 =\{\eta:(\vartheta,\eta)\in\Theta\}.
\]
于是 剖面 LR 统计量严格定义为
\begin{equation*}\operatorname{LR}_n^{\rm prof}(\vartheta)
 =2\left\{
 \ell_n(\widehat\theta_n)
 -\sup_{\eta\in\mathcal H(\vartheta)}\ell_n(\vartheta,\eta)
 \right\},
 \qquad \vartheta\in\mathcal V.
\end{equation*}
在本章前述正则条件下，真值 $\vartheta_0$ 满足
\[
 \operatorname{LR}_n^{\rm prof}(\vartheta_0)\xrightarrow{d}\chi^2_r.
\]
因此
\begin{equation*}C^{\rm PL}_{1-\alpha}
 =\left\{
 \vartheta\in\mathcal V:
 \operatorname{LR}_n^{\rm prof}(\vartheta)\le\chi^2_{r,1-\alpha}
 \right\}
\end{equation*}
具有渐近覆盖率 $1-\alpha$。它不是有限样本精确区间；边界参数、奇异信息、支撑依赖参数或不可识别模型中，临界分布必须重新求出。
\end{definition}

\begin{proposition}[剖面 LR 与有效信息椭球]
剖面 LR 区域与 Wald 区域的一阶等价可以直接从有效信息看出。前文消去 干扰参数 后的 剖面局部渐近正态 二次型为
\[
 q_n^{\rm prof}(h)
 =C_n+h^{\mathsf T}\Delta_{\vartheta\cdot\eta,n}
 -\frac12h^{\mathsf T}I_{\vartheta\cdot\eta}h
 +o_{\mathbb P}(1).
\]
其无约束顶点为
\[
 \widehat h
 =I_{\vartheta\cdot\eta}^{-1}
 \Delta_{\vartheta\cdot\eta,n}.
\]
把任意 $h$ 与顶点比较，完成平方得到
\begin{align*}
 2\{q_n^{\rm prof}(\widehat h)-q_n^{\rm prof}(h)\}
 &=(h-\widehat h)^{\mathsf T}
 I_{\vartheta\cdot\eta}
 (h-\widehat h)
 +o_{\mathbb P}(1).
\end{align*}
所以 剖面 LR 置信集在 $n^{-1/2}$ 局部尺度上正是 Fisher 有效信息诱导的椭球。Wald 区域是在估计点处直接写出同一个椭球；得分 区域则在候选原假设点处用有效 得分 测量法向距离。三者有限样本边界可以不同，但一阶几何相同。
\end{proposition}

本章的统一图像是：正则 似然 在 $n^{-1/2}$ 局部尺度上是带 Fisher 度量的随机二次型，检验只测量无约束极大点相对零假设切空间的法向距离；干扰参数 是先投影掉的切向方向，Pearson 检验则把完全相同的几何搬到概率单纯形。下一章中，线性模型会把这些抽象的 Fisher 投影变成观测空间中的显式矩阵投影，并在正态误差下给出有限样本精确 $F$ 理论。








